\clearpage
\begingroup
\tcbset{myboxstyle/.append style={breakable=false}}
\everymath{\textstyle}
\subsection{Счётная аддитивность продолжения меры}

Пусть $S$ --- полуалгебра подмножеств $\Omega$, то есть полукольцо,
содержащее $\Omega$, а $m$ --- конечная $\sigma$-аддитивная мера на $S$.
Обозначим через $\gamma$ её продолжение на $R(S)$:
\[
    \gamma\left(\bigsqcup_{i=1}^{n}A_i\right)
    =\sum_{i=1}^{n}m(A_i),\qquad A_i\in S.
\]
Поскольку $\Omega\in S$, кольцо $R(S)$ является алгеброй.

\begin{theorem}
    Продолжение $\gamma$ является $\sigma$-аддитивной мерой на $R(S)$.
    В частности, оно $\sigma$-субаддитивно.
\end{theorem}

\subsection{Внешняя мера Лебега}

\begin{definition}
    Для любого $A\subseteq\Omega$ положим
    \[
        \mu^*(A)=\inf\left\{
            \sum_{i=1}^{\infty}m(A_i):\
            A\subseteq\bigcup_{i=1}^{\infty}A_i,\quad A_i\in S
        \right\}.
    \]
    Функция $\mu^*\colon2^\Omega\to[0;+\infty)$ называется
    \textbf{внешней мерой Лебега}, порождённой $m$.
    Инфимум берётся по всем счётным покрытиям $A$ элементами $S$.
\end{definition}

\begin{remark}[Конечность внешней меры]
    Любое $A\subseteq\Omega$ покрывается множеством $\Omega\in S$, поэтому
    \[
        0\leq\mu^*(A)\leq m(\Omega)<\infty.
    \]
    Конечное покрытие можно дополнить пустыми множествами.
    В частности, $\mu^*(\varnothing)=0$.
\end{remark}

\begin{theorem}[Совпадение с мерой на алгебре]
    Для любого $A\in R(S)$ справедливо равенство $\mu^*(A)=\gamma(A)$.
\end{theorem}

\begin{Proof}
    Пусть $A=\bigsqcup_{i=1}^{n}A_i$, $A_i\in S$.
    Это покрытие $A$, поэтому
    \[
        \mu^*(A)\leq\sum_{i=1}^{n}m(A_i)=\gamma(A).
    \]
    С другой стороны, для любого покрытия $A\subseteq\bigcup_{j=1}^{\infty}B_j$,
    $B_j\in S$, по $\sigma$-субаддитивности $\gamma$
    \[
        \gamma(A)\leq\sum_{j=1}^{\infty}\gamma(B_j)
        =\sum_{j=1}^{\infty}m(B_j).
    \]
    Переходя к инфимуму по покрытиям, получаем $\gamma(A)\leq\mu^*(A)$.
\end{Proof}

\newpage
\subsection{Свойства внешней меры}

\begin{example}[Внешняя мера не обязана быть аддитивной]
    Пусть
    \[
        \Omega=\{1,2,3\},\qquad
        S=\{\varnothing,\{1,2\},\{3\},\Omega\},
    \]
    а мера на $S$ задана значениями
    \[
        m(\varnothing)=0,\quad m(\{1,2\})=2,\quad
        m(\{3\})=0,\quad m(\Omega)=2.
    \]
    Тогда $\mu^*(\{1\})=\mu^*(\{2\})=2$ и $\mu^*(\{1,2\})=2$.
    Следовательно,
    \[
        \mu^*(\{1\}\sqcup\{2\})=2
        \ne\mu^*(\{1\})+\mu^*(\{2\})=4.
    \]
\end{example}

\begin{theorem}[Счётная субаддитивность]
    Для любых $A,A_1,A_2,\ldots\subseteq\Omega$
    \[
        A\subseteq\bigcup_{i=1}^{\infty}A_i
        \quad\Longrightarrow\quad
        \mu^*(A)\leq\sum_{i=1}^{\infty}\mu^*(A_i).
    \]
\end{theorem}

\begin{Proof}
    Если сумма справа бесконечна, утверждение очевидно.
    Пусть она конечна. Зафиксируем $\eps>0$. По определению инфимума
    для каждого $i$ выберем покрытие
    \[
        A_i\subseteq\bigcup_{j=1}^{\infty}B_{ij},\qquad B_{ij}\in S,
        \qquad
        \sum_{j=1}^{\infty}m(B_{ij})<\mu^*(A_i)+\frac{\eps}{2^i}.
    \]
    Счётное семейство всех $B_{ij}$ покрывает $A$. Поэтому
    \[
        \begin{aligned}
            \mu^*(A)
            &\leq\sum_{i=1}^{\infty}\sum_{j=1}^{\infty}m(B_{ij})\\
            &\leq\sum_{i=1}^{\infty}\mu^*(A_i)
                 +\sum_{i=1}^{\infty}\frac{\eps}{2^i}
             =\sum_{i=1}^{\infty}\mu^*(A_i)+\eps.
        \end{aligned}
    \]
    Устремляя $\eps$ к нулю, получаем требуемое неравенство.
\end{Proof}

\newpage
\begin{theorem}[Оценка через симметрическую разность]
    Для любых $A,B\subseteq\Omega$
    \[
        |\mu^*(A)-\mu^*(B)|\leq\mu^*(A\mathbin{\triangle}B).
    \]
\end{theorem}

\begin{Proof}
    Из включений
    \[
        A\subseteq B\cup(A\mathbin{\triangle}B),\qquad
        B\subseteq A\cup(A\mathbin{\triangle}B)
    \]
    и субаддитивности следует
    \[
        \begin{aligned}
            \mu^*(A)&\leq\mu^*(B)+\mu^*(A\mathbin{\triangle}B),\\
            \mu^*(B)&\leq\mu^*(A)+\mu^*(A\mathbin{\triangle}B).
        \end{aligned}
    \]
    Все величины конечны, поэтому эти два неравенства дают требуемую оценку.
\end{Proof}

\subsection{Измеримые по Лебегу множества}

\begin{definition}
    Множество $A\subseteq\Omega$ называется \textbf{измеримым по Лебегу}
    относительно построенной внешней меры, если
    \[
        \forall\eps>0\quad\exists A_\eps\in R(S):\quad
        \mu^*(A\mathbin{\triangle}A_\eps)<\eps.
    \]
    Обозначим систему всех таких множеств через $M$.
\end{definition}

\begin{theorem}
    $R(S)\subseteq M$.
\end{theorem}

\begin{Proof}
    Для $A\in R(S)$ и любого $\eps>0$ выберем $A_\eps=A$.
    Тогда $\mu^*(A\mathbin{\triangle}A_\eps)=\mu^*(\varnothing)=0<\eps$.
\end{Proof}

\begin{theorem}
    Система $M$ измеримых по Лебегу множеств является алгеброй.
\end{theorem}

\newpage
\begin{Proof}[Замкнутость относительно пересечения]
    Пусть $A,B\in M$ и $\eps>0$. Выберем
    $A_{\eps/2},B_{\eps/2}\in R(S)$ так, чтобы
    \[
        \mu^*(A\mathbin{\triangle}A_{\eps/2})<\frac\eps2,
        \qquad
        \mu^*(B\mathbin{\triangle}B_{\eps/2})<\frac\eps2.
    \]
    Тогда $A_{\eps/2}\cap B_{\eps/2}\in R(S)$ и
    \[
        (A\cap B)\mathbin{\triangle}(A_{\eps/2}\cap B_{\eps/2})
        \subseteq
        (A\mathbin{\triangle}A_{\eps/2})
        \cup(B\mathbin{\triangle}B_{\eps/2}).
    \]
    По субаддитивности
    \[
        \begin{aligned}
            &\mu^*\bigl((A\cap B)\mathbin{\triangle}
                           (A_{\eps/2}\cap B_{\eps/2})\bigr)\\
            &\qquad\leq\mu^*(A\mathbin{\triangle}A_{\eps/2})
                      +\mu^*(B\mathbin{\triangle}B_{\eps/2})<\eps.
        \end{aligned}
    \]
    Следовательно, $A\cap B\in M$.
\end{Proof}

\begin{Proof}[Замкнутость относительно симметрической разности]
    Для тех же приближений
    $A_{\eps/2}\mathbin{\triangle}B_{\eps/2}\in R(S)$ и
    \[
        \begin{aligned}
            &(A\mathbin{\triangle}B)\mathbin{\triangle}
                (A_{\eps/2}\mathbin{\triangle}B_{\eps/2})\\
            &\qquad\subseteq
                (A\mathbin{\triangle}A_{\eps/2})
                \cup(B\mathbin{\triangle}B_{\eps/2}).
        \end{aligned}
    \]
    Поэтому
    \[
        \begin{aligned}
            &\mu^*\bigl((A\mathbin{\triangle}B)\mathbin{\triangle}
                (A_{\eps/2}\mathbin{\triangle}B_{\eps/2})\bigr)\\
            &\qquad\leq\mu^*(A\mathbin{\triangle}A_{\eps/2})
                      +\mu^*(B\mathbin{\triangle}B_{\eps/2})<\eps.
        \end{aligned}
    \]
    Следовательно, $A\mathbin{\triangle}B\in M$.
    Замкнутость относительно пересечения и симметрической разности
    означает, что $M$ --- кольцо. Так как $\Omega\in R(S)\subseteq M$,
    это кольцо является алгеброй.
\end{Proof}
\endgroup
